% !TeX encoding = UTF-8 !TeX spellcheck = en_US
\documentclass[11pt, a4paper]{article}

% The title of the current document to be produced.
\newcommand{\doctitle}{Course Outline}
\newcommand{\assignments}{\bluetext{\textbf{Assignment:}}}
%
\setlength{\unitlength}{1in}
\renewcommand{\arraystretch}{2}
\setlength{\columnsep}{2cm}

\usepackage[sfdefault]{atkinson}
\usepackage[T1]{fontenc}


%------------------------------------------------------------
% Import commands for both teacher and course information.  | NOTE: Change your
% teacher and course info in these files. |
%------>------>------>------>------>------>------>------>-->|
%% ==================================
%% ===== Teacher info             ===
%% ==================================
%%
\newcommand{\instructor}{Katie Breivik}
\newcommand{\office}{Wean 8319}
\newcommand{\hours}{Tuesday at 3pm or by email appointment}
\newcommand{\phone}{}
\newcommand{\college}{}
\newcommand{\email}{kbreivik@andrew.cmu.edu}
\newcommand{\faculty}{}
\newcommand{\department}{Department of Physics}
                              %|
%% ==================================
%% ===== Course-specific commands ===
%% ==================================

%- Instructions: change course info here. 
\newcommand{\semester}{Spring 2025}

\newcommand{\ponderation}{2-4-3 (Theory-Lab-Homework)}
\newcommand{\coursetitle}{Science and Science Fiction}
\newcommand{\coursenumber}{[33-120]}
\newcommand{\prerequisite}{None}
                               %|   
%
%------------------------------------------------------------
%-- Import packages and custom command definitons.          |
%------>------>------>------>------>------>------>------>-->|
\input{includes/packages-imports}                          %|  
\input{includes/custom-commands}   
%
%---> Genereate & inject metadata                           |
%--------------------------------------------------------------
%-- Generate and inject metadate in the produced PDF document |
%------>------>------>------>------>------>------>------>-->---
\hypersetup{pdfauthor={\instructor},%
	pdftitle={\coursenumber -- \coursetitle},%
	pdfkeywords={\college,  \department},%
	pdfproducer={LaTeX},%
	pdfcreator={pdfLaTeX},
	bookmarks,
	bookmarksnumbered = true,
	bookmarksopen     = true,
	pdfpagelabels     = true,
	pdfstartview={XYZ null null 1.2}
}                          %|
%------------------------------------------------------------

\topmargin -70pt
\begin{document} 

\normalfont

%-------------------------------------------------------------
%-- Make the header of the document                          |
%------>------>------>------>------>------>------>------>--> |
\input{includes/document-header}
  
%-------------------------------------------------------------
%-- Insert the course & teacher info                         |
%------>------>------>------>------>------>------>------>--> |
\hrule     
\vspace{.5cm}

\begin{multicols}{2}
    \begin{description}[labelindent=0.00in,leftmargin=1.25in,rightmargin=0.0in,style=nextline]
        %--> First column:         
		\item[\textsc{Instructor}:] \instructor
        \item[\textsc{E-mail}:] \email
        \item[] 
        \item[\textsc{Class Time}:] Mon/Wed/Fri\newline   11-11:50am 
        \item[] 
        %--> Second column:         
        \item[\textsc{Office}:]  \office
        \item[\textsc{Office Hours}:] \hours
        \item[] 
        \item[\textsc{Class Room}:] DH A302
        \item[] 
    \end{description}
\end{multicols}
\hrule        
\vspace{.2cm}

 %--------  Course Description  ------------------------------
\customsection{Course Description from Course Catalog}  
\noindent We will view and critique the science content in a selection of
science fiction films, spanning more than 100 years of cinematic history, and
from sci-fi TV shows from the past 50+ years. Guided by selected readings from
current scientific literature, and aided by order-of-magnitude estimates and
careful calculations, we will ponder whether the films are showing things which
may fall into one of the following categories: Science fiction at the time of
production, but currently possible, due to recent breakthroughs. Possible, in
principle, but beyond our current technology. Impossible by any science we know.
Topics to be covered include the future of the technological society, the
physics of Star Trek, the nature of space and time, extraterrestrial
intelligence, robotics and artificial intelligence, biotechnology and more.
Success of this course will depend upon class participation. Students will be
expected to contribute to discussion of assigned readings and problems.

\customsection{Student Learning Outcomes} 
\noindent At the end of the class, students should be able to
\begin{borderedsquare}
     \setlength\itemsep{0.3em}        
	\item develop and maintain an interest in science.
	\item become familiar with the history of science fiction film by watching
    numerous short clips and full-length movies.
    \item develop critical thinking skills and to discern the difference between
    currently plausible scenarios, future possibilities, and total fantasy.
    \item use science fiction film as a springboard for discussing current
    scientific research.
    \item use library and/or online resources and to discern and summarize the
    key ideas in short scientific articles.
    \item work cooperatively in teams to solve challenging, but doable problems.
\end{borderedsquare}

\customsection{Textbooks and Course Materials}  
%---------------------------------
%--> List of recommended textbooks. 
\begin{itemize}[itemsep=4pt,parsep=0pt,topsep=1pt,partopsep=1pt]
	\item[\color{darkblue}\faNewspaperO]\textbf{Canvas:} The site can be reached
	from \href{http://www.cmu.edu/canvas/}{http://www.cmu.edu/canvas/} The
	course website will contain the syllabus, homework, project assignments and
	all required reading. Announcements and visual material presented during
	lectures will also be posted throughout the semester.\\

	\item[\color{darkblue}\faBook] \textbf{\textsc{Textbooks:}} Exploring
    Science Through Science Fiction, \textbf{Second Edition}, Author: Barry B.
    Loukkala, available on
    \href{https://www.amazon.com/Exploring-Science-Through-Fiction-dp-3030293920/dp/3030293920/ref=dp_ob_title_bk}{Amazon}
\end{itemize}

\customsection{Prerequisite Knowledge: No prequisite knowledge assumed}

\customsection{Evaluation Procedures}    
\label{sec:eval-procudures}
\noindent The course will consist of in-class lectures and a variety of
assignments, team projects, and quizzes as detailed below. There will be a final
cumulative exam to be taken online on Canvas during finals week. 

\begin{itemize}[itemsep=8pt,parsep=0pt,topsep=3pt,partopsep=4pt]
    \item[]
%    \begin{center}
        \renewcommand{\arraystretch}{1.5} % this reduces the vertical spacing
        %between rows    
        \begin{tabular}{lll}
            \thead{\color{darkblue} Course Component} & \thead{\color{darkblue}
            Breakdown}& \thead{\color{darkblue} Total Points} \\ 
            %---s Load the table body: dynamic table content.        
       \hline    
            \textbf{Science Knowledge Survey} & 1 on 1st day of class & 25 pts
            \\                         
        \hline    
	        \textbf{In-Class Discussion} & Every class period & 50 pts \\                         
        \hline    
	        \textbf{Sci-fi Problem Sets} & 4 at 25 pts each & 100 pts \\                         
        \hline    
	        \textbf{Exploration Papers} & 5 at 25 pts each & 125 pts \\                         
        \hline    
	        \textbf{Quizzes on Lectures/Reading} & 7 w/ variable pts & 350 pts
	        \\                         
        \hline
            \textbf{Team Projects} & 2 at 25 pts each & 50 pts \\                         
        \hline
            \textbf{Film Critiques} & 2 at 50 pts each & 100 pts \\                         
        \hline
            \textbf{Final Exam} &  & 200 pts \\                         
        \hline
        \hline
            \textbf{Total Points} & & 1000 pts\\
        



        \end{tabular}
%    \end{center} 
    %        \evaluationrules
\end{itemize} 

\noindent Final grade boundaries are: A: 900 - 1000, B: 800 - 899, C: 700 - 799
pts D: 600 - 699, R: <600 

\noindent Boundary cases may be discussed based on demonstrated engagment,
effort, and improvement over the semester. 

\vspace{.4cm}

\customsection{Attendance and Late Work Policies}  
\noindent 
Formal attendance will not be recorded for the course. However, to gain
discussion points for the course, you must be present in class. Class time will
also serve as a venue to communicate important details on course component, so
in-class participation is likely to correlate strongly with final grade
outcomes.

Late homeworks may be submitted for reduced credit at 2 pts per day up to the
point that grading is complete. Assignments are allowed \emph{one submission
only} and there are \emph{\textbf{no makeup quizzes.}} The dates of each quiz
are listed in the course calendar and will be announced both in class and on
canvas. There will be \emph{several} opportunities for extra credit announced
during class that can be used to make up points lost for a missed quiz. 

\vspace{.4cm}

\customsection{Accommodations for Students with Disabilities}
\noindent
If you have an accommodations letter from the Disability Resources office, you
are encouraged to discuss your accommodations and needs with me as early in the
semester as possible. I will work with you to ensure that accommodations are
provided as appropriate. If you suspect that you may benefit from accommodations
but are not yet registered with the Office of Disability Resources, I encourage
you to contact them at access@andrew.cmu.edu.
\vspace{.4cm}


\customsection{Statement on Diversity, Equity, and Inclusion}  
\noindent 
Everyone in this course must be treated with respect. Science is a human
endeavour and we are all accountable for ensuring that the environment we create
is welcoming and free from discrimination. Each of us is here for a variety of
reasons, comes from a variety of backgrounds, and have a diversity of skills and
past experiences. It is expected that the students and instructor will
contribute to a collaborative learning environment where diverse perspectives
and approaches are celebrated.

Unfortunately, incidents may occur, whether intentional or not, that can create
an unwelcoming environment for one or more class members. The university
encourages anyone who experiences or observes unfair or hostile treatment on the
basis of identity to speak out for justice and support, within the moment of the
incident or after the incident has passed. Anyone can share these experiences
using the following resources:
\begin{itemize}
    \item Center for Student Diversity and Inclusion: csdi@andrew.cmu.edu, (412)
268-2150 
    \item Report-It online anonymous reporting platform: reportit.net username:
tartans password: plaid 
\end{itemize}
\noindent All reports will be documented and deliberated to determine if there
should be any following actions. Regardless of incident type, the university
will use all shared experiences to transform our campus climate to be more
equitable and just.

\vspace{.4cm}

\customsection{Communication Expectations} 
\noindent The majority of our communication will be done during class time but I
am also reachable by email at the address listed at the top of this page. I will
also hold a weekly office hour on Tuesdays at 2pm. I can also meet outside of
those hours, but these meetings will need to scheduled by email. 
\vspace{.4cm}

\customsection{Use of Film/TV Clips Under Fair Use} 
\noindent We will view several clips taken from television and film in class for
illustrative and critical purposes in support of the aims of the course. They
will only be shown during class and not available for personal use such that all 
clips remain under fair use.   
\vspace{.4cm}

\customsection{Student Wellness}
\noindent Take care of yourself. Do your best to maintain a healthy lifestyle
this semester by eating well, exercising, getting enough sleep and taking some
time to relax. This will help you achieve your goals and cope with stress. 

All of us benefit from support during times of struggle. You are not alone.
There are many helpful resources available on campus and an important part of
the college experience is learning how to ask for help. Asking for support
sooner rather than later is often helpful. If you or anyone you know experiences
any academic stress, difficult life events, or feelings like anxiety or
depression, we strongly encourage you to seek support including from the
\href{https://www.cmu.edu/wellbeing/resources/timely-care.html}{Timely Care
App}. Counseling and Psychological Services (CaPS) is alsoå here to help: call
412-268-2922 and visit their \href{http://www.cmu.edu/counseling/}{website}.

Consider reaching out to a friend, faculty or family member you trust for help
getting connected to the support that can help. If you or someone you know is
feeling suicidal or in danger of self-harm, call someone immediately, day or
night: CaPS: 412-268-2922 Re:solve Crisis Network: 888-796-8226

The \href{https://www.cmu.edu/student-affairs/resources/cmu-pantry/}{CMU Pantry}
is committed to reducing hunger among students by providing nutritious food at
no cost. Accessing these food resources can promote a healthier, balanced
university experience.

\clearpage

\end{document} 